\chapter{Bài toán và thiết kế đánh giá}\label{chap:problem}
\FloatBarrier
\section{Bối cảnh và phát biểu bài toán}
Mỗi bài báo VietNamNet được tổ chức theo một chuyên mục biên tập, nhưng nội dung có thể giao thoa nhiều lĩnh vực. Bài toán đặt ra là dự đoán một chuyên mục chính từ tiêu đề và nội dung, trong khi nhãn tham chiếu phản ánh taxonomy của tòa soạn thay vì một hệ phân loại chủ đề phổ quát. Các khó khăn gồm phân đoạn từ tiếng Việt, độ dài bài báo, chênh lệch số mẫu giữa các lớp và ranh giới nhãn không hoàn toàn tách biệt. Chất lượng nguồn nhãn và tính nhất quán của phép chia dữ liệu vì vậy có vai trò quan trọng bên cạnh lựa chọn bộ phân loại.

Với $x_i=(t_i,b_i)$ gồm tiêu đề và nội dung, $y_i\in\mathcal C$, $|\mathcal C|=20$, corpus giám sát là
\begin{equation}\mathcal D=\{(x_i,y_i)\}_{i=1}^{199\,965}.\end{equation}
Phương pháp $m$ biến đổi văn bản bằng $\phi_m$, sinh điểm với mô hình $f_{\theta_m}$ và chọn nhãn:
\begin{equation}
z_i^{(m)}=f_{\theta_m}(\phi_m(x_i)),\qquad
\hat y_i^{(m)}=\arg\max_{c\in\mathcal C}z_{i,c}^{(m)}.
\end{equation}
Argmax là giao thức đối chứng chính. Riêng PhoBERT có thêm quy tắc quyết định theo lớp để phục vụ ứng dụng, được đánh giá tách biệt. URL và breadcrumb dùng xác minh nguồn nhãn, không trở thành đặc trưng của bộ phân loại. Top-$k$ là các ứng viên cạnh tranh, không phải một tập nhãn cùng đúng: tác vụ là đa lớp đơn nhãn.

Một bài về sinh kế vùng dân tộc có thể đồng thời chứa du lịch, văn hóa và chính sách. Vì thế, sai số được hiểu là không khớp chuyên mục tham chiếu, chưa mặc nhiên là sai theo mọi cách diễn giải nội dung. Mô hình không kiểm chứng thông tin trong bài và không có bộ phát hiện chắc chắn văn bản ngoài taxonomy.

\FloatBarrier
\section{Mục tiêu nghiên cứu và phạm vi}
SVM, Logistic Regression, BiLSTM và PhoBERT được lựa chọn để khảo sát ba họ biểu diễn: đặc trưng từ vựng, embedding chuỗi và biểu diễn ngữ cảnh tiền huấn luyện. Báo cáo trả lời ba câu hỏi:
\begin{enumerate}[label=\textbf{RQ\arabic*:},leftmargin=17mm]
\item Trên cùng corpus VietNamNet và taxonomy 20 chuyên mục, chất lượng phân loại của bốn phương pháp khác nhau như thế nào?
\item Khác biệt chất lượng có tương xứng với thời gian huấn luyện, bộ nhớ và chi phí suy luận của từng hướng tiếp cận không?
\item Các phương pháp khác nhau thế nào ở chuyên mục ít mẫu hoặc có nội dung giao thoa?
\end{enumerate}
Notebook dữ liệu tạo đầu vào chung; bốn notebook huấn luyện tự chứa quy trình tạo biểu diễn, chọn cấu hình, đánh giá và xuất artifact. LSTM được thực thi bằng biến thể hai chiều BiLSTM, với embedding học từ đầu. Các gói đối chứng phục vụ đánh giá offline; ứng dụng minh họa dùng PhoBERT để đề xuất nhãn, hỗ trợ rà soát và lưu phản hồi (Phụ lục~\ref{app:application}).

So sánh xem xét hiệu quả toàn bộ pipeline trong một điều kiện thực nghiệm cụ thể. Nó chưa tách riêng đóng góp của kiến trúc, pretraining, stopword, lặp tiêu đề hay vùng văn bản được giữ. Phân tích lớp hiếm và độ tin cậy xác suất bổ sung cho điểm tổng, nhưng không thay thế đánh giá ngoài nguồn hoặc nhiều lần train.


Nhãn cuối là một quyết định của bài toán đơn nhãn, còn confidence phục vụ đánh giá mức chắc chắn trong điều kiện taxonomy. Một bài ngoài miền vẫn có thể được gán nhãn cao điểm do bộ phân loại chỉ có 20 lựa chọn; không coi tổng xác suất bằng 1 như bằng chứng đầu vào hợp lệ. Vì thế, ba câu hỏi nghiên cứu không chỉ hỏi model nào có điểm cao nhất, mà còn hỏi lỗi rơi vào đâu và số đo nào có thể hỗ trợ sử dụng kết quả.

Phạm vi ứng dụng là vòng rà soát có người quyết định, không gồm tự sinh bài hoặc tự cập nhật trọng số từ feedback. Bằng chứng giao diện giải thích cách kết quả NLP được dùng; nó không được cộng vào test hoặc dùng chứng minh tổng quát hóa. Sự tách biệt giữa thực nghiệm và minh họa giúp cùng một chỉ số không bị gán hai ý nghĩa khác nhau.

\FloatBarrier
\section{Giao thức dữ liệu và tiêu chí đánh giá}
Khử trùng diễn ra trước chia tập. Phép chia phân tầng 70/15/15, seed 42 tạo train 139.975 bài và hai tập validation, test, mỗi tập 29.995 bài. Bốn pipeline có cùng chỉ số chia, URL và nhãn test theo thứ tự. Train dùng cập nhật trọng số, tính class weights và fit vocabulary/IDF; validation chọn cấu hình, checkpoint và tham số hiệu chỉnh; test chỉ đánh giá cuối sau khi các lựa chọn đã cố định.
\SplitDiagram
Macro-F1 là tiêu chí lựa chọn để mỗi lớp đóng góp ngang nhau; Accuracy và Weighted-F1 bổ sung góc nhìn theo số bài. Kết quả theo lớp, confusion matrix và lỗi truy vết chỉ ra nơi chất lượng suy giảm. Chương~\ref{chap:experiments} đánh giá riêng chất lượng nhãn, xác suất và chi phí, tránh lấy một chỉ số thay thế tất cả yêu cầu.

Validation được dùng cho cả lựa chọn model và hiệu chỉnh. Sự dùng lại này tạo phụ thuộc vào tập validation, dù test không tham gia fit. Một seed hỗ trợ tái tạo phép chia trong snapshot, chưa đo biến thiên huấn luyện. Chín đối chiếu giữa các tập theo URL, văn bản thô và văn bản chuẩn hóa không phát hiện giao nhau chính xác; phạm vi kiểm tra chưa bao gồm gần trùng hoặc bài cùng sự kiện.

Bốn pipeline được đánh giá trên cùng corpus sau khử trùng và cùng tập test; các phép học biểu diễn riêng chỉ dùng train. Chương tiếp theo xem xét bằng chứng nghiên cứu làm cơ sở lựa chọn các phương pháp.
